\documentclass[12pt]{article}
\usepackage[utf8]{inputenc}
\usepackage{graphicx}
\usepackage{amsmath}
\usepackage{geometry}
\usepackage{array}
\usepackage{longtable}
\usepackage{booktabs}
\geometry{letterpaper, margin=1in}

\begin{document}

\begin{flushright}
Gabriela Mora, V-31.541.893 \\
Lucía Martínez, V-31.906.577 \\
\end{flushright}

\vspace{1cm}

\begin{center}
\LARGE\textbf{Práctica del Laboratorio 1} \\
\vspace{0.3cm}
\large Arquitectura del Computador \\
\vspace{0.3cm}
\normalsize 29 de Julio de 2026
\end{center}

\vspace{1cm}

\section*{1. ¿Cómo se implementa la recursividad en MIPS32?}

La recursividad en MIPS32 se implementa utilizando la pila (stack) para preservar el estado de cada llamada recursiva. El registro \$sp (stack pointer) es fundamental en este proceso.

\subsection*{¿Qué papel cumple la pila (\$sp)?}

El stack pointer (\$sp) cumple estas funciones clave:

\begin{itemize}
    \item \textbf{Almacenamiento de direcciones de retorno:} Guarda \$ra para saber dónde retornar después de cada llamada.
    \item \textbf{Preservación de registros:} Guarda valores de registros que podrían modificarse en llamadas posteriores.
    \item \textbf{Almacenamiento de parámetros y variables locales:} Para cada instancia de la llamada recursiva.
\end{itemize}

\section*{2. ¿Qué riesgos de desbordamiento existen?}

\subsection*{1. Desbordamiento Aritmético (Arithmetic Overflow)}

Ocurre cuando el resultado de una operación supera los 32 bits, provocando errores en el bit de signo o pérdida de datos.

\subsubsection*{Cómo mitigarlo (Ejemplo en MIPS):}

En lugar de utilizar \texttt{add \$t0, \$t1, \$t2} (que dispara excepción), emplear \texttt{addu \$t0, \$t1, \$t2}. Esta instrucción suma los bits sin verificar el acarreo hacia el bit de signo, evitando interrupciones inesperadas. Para validar manualmente, se puede usar \texttt{slt \$t3, \$t1, \$zero} y comparar con el resultado, asegurando que el programa maneje el error mediante un salto lógico antes de guardar valores corruptos.

\subsection*{2. Desbordamiento de Estructura (Stack Overflow)}

Se origina cuando la recursión excesiva agota la memoria reservada para el puntero de pila (\$sp), invadiendo segmentos de memoria ajenos.

\subsubsection*{Cómo mitigarlo (Ejemplo en MIPS):}

\begin{itemize}
    \item \textbf{Caso base:} Usar \texttt{beq \$a0, \$zero, finparidad} para detener la recursión inmediatamente.
    \item \textbf{Validación de rango:} Emplear \texttt{bltz \$a0, fin} al inicio para evitar que números negativos provoquen una recursión infinita.
    \item \textbf{Contador de seguridad:} Se puede utilizar \texttt{slti \$t0, \$s1, 100} (donde \$s1 es un contador de profundidad) antes de ejecutar el \texttt{jal paridad}. Si el contador llega a 100, el programa salta a una rutina de error (\texttt{j error}) antes de que el \$sp corrompa la memoria.
\end{itemize}

\subsection*{3. Desbordamiento de Tipo (Type Overflow / Truncation)}

Ocurre al intentar comprimir datos de 64 bits (procedentes de \texttt{mult}) en registros de 32 bits, o por el uso de direcciones mal alineadas.

\subsubsection*{Cómo mitigarlo en MIPS 32:}

\begin{itemize}
    \item Extraer los bits superiores con \texttt{mfhi} y los bits inferiores con \texttt{mflo} después de una multiplicación.
    \item Verificar que el valor en \texttt{mfhi} sea igual a la extensión de signo del registro LO.
    \item Guardar ambos registros (HI y LO) por separado en memoria si se necesita preservar el resultado completo.
    \item Asegurar que los offsets sean siempre múltiplos de 4 al usar \texttt{lw} o \texttt{sw} para evitar errores de alineación.
\end{itemize}

\section*{3. ¿Qué diferencias encontraste entre una implementación iterativa y una recursiva en cuanto al uso de memoria y registros?}

\subsection*{1. Uso de Memoria}

\textbf{Iterativa:} Es mucho más eficiente en cuanto a memoria. El código iterativo no necesita la pila. Todo el trabajo se realiza usando unos pocos registros (como \$t0, \$t1) que se sobrescriben en cada vuelta del bucle. No importa si el número es pequeño o grande; el consumo de memoria es siempre constante ($O(1)$).

\textbf{Recursiva:} Es costosa. Cada vez que la función se llama a sí misma, el programa debe reservar espacio en la pila (\$sp) para guardar la dirección de retorno (\$ra) y los argumentos. Si el número a evaluar es grande, la pila crece linealmente ($O(n)$), lo que consume RAM rápidamente.

\subsection*{2. Uso de Registros}

\textbf{Iterativa:} Los registros son reutilizados. Una vez que el bucle pasa a la siguiente iteración, el valor anterior ya no es necesario, por lo que se reemplaza. Es un proceso "destructivo" pero muy rápido.

\textbf{Recursiva:} Los registros deben ser "preservados". Antes de llamar a la siguiente instancia de la función, el código debe copiar el valor actual de los registros a la memoria (pila) para "congelar" el estado. Al volver de la recursión, debe restaurarlos (\texttt{lw}). El procesador gasta más tiempo moviendo datos entre registros y memoria que haciendo la operación matemática en sí.

\section*{4. ¿Qué diferencias encontraste entre los ejemplos académicos del libro y un ejercicio completo y operativo en MIPS32?}

Al comparar los ejemplos académicos con nuestra implementación real, identificamos tres diferencias clave:

\begin{enumerate}
    \item \textbf{El entorno obligatorio:} A diferencia de los fragmentos de código del libro, nuestro programa requirió directivas obligatorias (\texttt{.data}, \texttt{.text}, \texttt{.globl main}) para que el simulador pudiera asignar memoria y ejecutar la lógica.
    \item \textbf{La optimización de registros:} Mientras el libro prioriza el uso de registros \$s (que requieren respaldo en pila), nosotros utilizamos registros \$t en la versión iterativa. Esto eliminó instrucciones \texttt{sw/lw} innecesarias, haciendo el código más ágil y eficiente.
    \item \textbf{Interacción con el usuario:} Los ejemplos del libro se limitan a cálculos matemáticos puros. En la práctica, tuvimos que integrar \texttt{syscalls} (servicios del sistema) para leer entradas y mostrar resultados, gestionando registros adicionales como \$v0 y \$a0.
\end{enumerate}

\section*{5. Elaborar un tutorial de la ejecución paso a paso en MARS.}

Este procedimiento aplica para el análisis de nuestros programas de paridad (iterativa y recursiva):

\begin{enumerate}
    \item \textbf{Ensamblado:} Abra el archivo .asm en MARS, acceda al menú Run y seleccione Assemble (F3). Esto prepara el código para su ejecución.
    \item \textbf{Preparación del entorno:} Asegúrese de tener visibles las pestañas Registers (para monitorear cambios en \$t, \$a y \$v) y Data Segment (esencial para observar el comportamiento de la pila \$sp en el ejercicio recursivo).
    \item \textbf{Depuración (Single Step):} Utilice el botón Single Step (F7) para ejecutar una instrucción a la vez.
    \begin{itemize}
        \item En el iterativo: Verifique cómo el registro \$t3 almacena el residuo tras la operación \texttt{div}.
        \item En el recursivo: Observe cómo el valor de \$sp decrece en la ventana Data Segment con cada llamada (\texttt{jal}), indicando la reserva de memoria, y cómo se recuperan los datos al ejecutar \texttt{lw}.
    \end{itemize}
    \item \textbf{Validación de resultados:} Una vez alcanzada la instrucción \texttt{syscall}, observe la consola Run I/O para confirmar que el resultado mostrado coincide con el valor calculado en los registros.
    \item \textbf{Reset:} Utilice el botón Reset para limpiar el estado y realizar nuevas pruebas con distintos números, garantizando que el programa sea consistente en cualquier escenario.
\end{enumerate}

\section*{6. Justificar la elección del enfoque (iterativo o recursivo) según eficiencia y claridad en MIPS.}

Tras implementar ambos métodos, elegimos el enfoque iterativo como el estándar óptimo por dos razones principales:

\begin{enumerate}
    \item \textbf{Eficiencia técnica:} La versión iterativa opera exclusivamente mediante registros, eliminando el alto costo de memoria (escritura/lectura en RAM) que exige la gestión de la pila en la recursión. Esto reduce significativamente los ciclos de reloj.
    \item \textbf{Claridad y seguridad:} Aunque la recursión es matemáticamente elegante, su gestión manual en ensamblador es propensa a errores. El enfoque iterativo es más directo, fácil de seguir paso a paso y elimina por completo el riesgo de un Stack Overflow, garantizando un funcionamiento robusto sin importar el tamaño del número ingresado.
\end{enumerate}

\section*{7. Análisis y Discusión de los Resultados}

\subsection*{Resultados cuantitativos}

Para un valor n cualquiera dentro del rango de 32 bits, observamos la siguiente diferencia en la ejecución:

\begin{table}[h]
\centering
\caption{Resultados cuantitativos (Paridad)}
\begin{tabular}{|c|c|c|}
\hline
\textbf{n} & \textbf{Ciclos (Iterativo)} & \textbf{Ciclos (Recursivo)} \\ \hline
10 & $\sim$20 & $\sim$150 \\ \hline
20 & $\sim$20 & $\sim$300 \\ \hline
30 & $\sim$20 & $\sim$450 \\ \hline
\end{tabular}
\end{table}

\vspace{0.5cm}

\begin{table}[h]
\centering
\caption{Comparativa técnica de enfoques}
\begin{tabular}{|l|l|l|}
\hline
\textbf{Métrica} & \textbf{Paridad Iterativa} & \textbf{Paridad Recursiva} \\ \hline
Complejidad de Tiempo & $O(1)$ & $O(n)$ \\ \hline
Consumo de Ciclos & Constante & Lineal \\ \hline
Uso de Memoria (Stack) & Nulo & $O(n)$ \\ \hline
\end{tabular}
\end{table}

\subsection*{Hallazgos principales}

\textbf{Límites arquitecturales:} Ambos métodos están limitados por la arquitectura de 32 bits de MIPS. Si el valor ingresado excede el rango máximo del registro ($2^{31} - 1$), se producirá un overflow aritmético en ambos casos. Además, en la versión recursiva, se presenta un límite adicional de Stack Overflow cuando n es lo suficientemente grande como para agotar el espacio reservado para la pila en el simulador.

\textbf{Optimizaciones efectivas:}

\begin{itemize}
    \item En el iterativo: Implementamos el uso directo de registros temporales (\$t), eliminando el acceso a memoria RAM y acelerando el cálculo.
    \item En el recursivo: Se minimizó el uso de registros guardados (\$s), optimizando el respaldo en la pila solo a lo estrictamente necesario (\$ra y \$a0).
\end{itemize}

\subsection*{Conclusiones}

La implementación iterativa es superior a la recursiva por los siguientes motivos:

\begin{enumerate}
    \item \textbf{Eficiencia computacional:} Al ser un proceso de tiempo constante ($O(1)$), no depende del tamaño de n, ejecutándose mucho más rápido que la recursiva, que debe realizar una llamada por cada unidad de n.
    \item \textbf{Optimización de recursos:} El enfoque iterativo opera exclusivamente en registros. Esto evita el consumo de memoria en la pila, lo cual es crítico para la estabilidad del sistema.
    \item \textbf{Estabilidad:} La versión iterativa es más robusta frente a la profundidad de llamadas. Mientras que la recursión está sujeta a fallos por desbordamiento de memoria de pila incluso antes de alcanzar los límites aritméticos del procesador, el método iterativo aprovecha al máximo el espacio de memoria disponible en la arquitectura.
\end{enumerate}

\end{document}
